\chapter{Multilingual Documents}

OmniLaTeX's translation infrastructure enables documents that mix multiple
languages within a single compilation.  This chapter provides practical
guidance for authoring multilingual content.

% =========================================================================
\section{Multi-Language Example}
% =========================================================================

The \texttt{examples/multi-language/} directory demonstrates a document
that combines English and German content.  The document is compiled with
English as the primary language and uses inline language switching for
German passages:

\begin{verbatim}
\documentclass[language=english]{omnilatex}
\begin{document}

\begin{english}
  This is the English abstract.
\end{english}

\begin{german}
  Dies ist die deutsche Zusammenfassung.
\end{german}

\end{document}
\end{verbatim}

Within each language environment, \cs{GetTranslation} resolves to the
appropriate localized string, and polyglossia applies the correct
hyphenation patterns.  Captions for figures and tables also switch
automatically when inside a language environment.

The example includes a full bibliography and glossary setup that
demonstrates how these subsystems interact with language switching.
See \texttt{examples/multi-language/README.md} for build instructions.

% =========================================================================
\section{Bibliography in Multiple Languages}
% =========================================================================

OmniLaTeX uses \texttt{biblatex} with the \texttt{ext-authoryear} style,
which supports multilingual citations out of the box.  Bibliography
headings are localized through the translation system:

\begin{verbatim}
\GetTranslation{FurtherReadingTitle}  % → "Further Reading" / "Weitere Literatur"
\end{verbatim}

When a bibliography entry contains titles in multiple languages,
\texttt{biblatex}'s \texttt{langid} field can be used to tag individual
entries:

\begin{verbatim}
@book{mueller2020,
  author    = {Müller, Hans},
  title     = {Deutsche Fachsprache},
  langid    = {german},
}
\end{verbatim}

The citation style automatically formats the entry according to its
language tag, including correct quotation marks and date formatting.

% =========================================================================
\section{Glossary in Multiple Languages}
% =========================================================================

The \texttt{glossaries-extra} package used by OmniLaTeX supports
multi-language labels through the translation system.  Glossary group
titles are resolved via \cs{GetTranslation}:

\begin{verbatim}
\glsxtrsetgrouptitle{greek}{\GetTranslation{Greek}}
\glsxtrsetgrouptitle{roman}{\GetTranslation{Roman}}
\end{verbatim}

This means glossary headings automatically appear as
``Griech.'' in German, ``Grec'' in French, or ``Greek'' in English,
depending on the active language context.

\cs{printunsrtglossary} can be called inside a language environment to
produce a glossary section with headings in the corresponding language.
For documents that maintain separate glossaries per language, call
\cs{printunsrtglossary} multiple times within different language
environments.

% =========================================================================
\section{Practical Tips}
% =========================================================================

\begin{itemize}
  \item \textbf{Consistent language use.}  Define a clear boundary between
    languages in your document.  Using language environments around entire
    sections or chapters is cleaner than switching mid-paragraph.

  \item \textbf{Translation completeness.}  OmniLaTeX ships translations
    for 47 keys across 18 languages (verified in
    \texttt{I18nCompleteness.lean}).  If you add custom keys, declare
    translations for all required languages to avoid silent English
    fallbacks.

  \item \textbf{Test all variants.}  Compile the document with each
    target language as the primary language (\texttt{language=german},
    \texttt{language=french}, etc.) to verify that captions, glossary
    headings, and translated strings render correctly.

  \item \textbf{Institution-specific keys.}  Institution modules may
    add their own translation keys (e.g.\ \texttt{TUHH},
    \texttt{LinkTUHH}).  Ensure these are also translated if you
    need multilingual output.
\end{itemize}
